\begin{tikzpicture}

\begin{axis}[
    width=\linewidth*\getDarcImageFactor,
    height=0.9*\linewidth*\getDarcImageFactor,
    xmin=-50,
    xmax=15,
    ymin=-90,
    ymax=38,
    xlabel={Eingangsleistung $P_\text{in}$},
    ylabel={Ausgangsleistung $P_\text{out}$},
    xtick={-50,-40,-30,-20,-10,0,10},
    ytick={-80,-60,-40,-20,0,20,30},
    yticklabels={},
    xticklabels={},
    grid=major,
    axis lines=left,
    axis line style={-Triangle},
    clip=false,
    legend style={draw, fill=white, font=\footnotesize, at={(0.98,0.02)}, anchor=south east},
]

\def\G{20}
\def\IIP{10}
\def\OIP{30}
\def\PSAT{21}

% ------------------------------------------------------------
% Sättigungsleistung
% ------------------------------------------------------------
\addplot[
    orange!90!black,
    very thick
] coordinates {
    (-50,\PSAT)
    (15,\PSAT)
};
\addlegendentry{Sättigung}

% ------------------------------------------------------------
% Reales Grundsignal mit Kompression und Sättigung
% ------------------------------------------------------------
\addplot[
    DARCblue,
    very thick,
    smooth
] coordinates {
    (-50,-30)
    (-40,-20)
    (-30,-10)
    (-20,0)
    (-10,9.5)
    (-5,14.0)
    (0,17.5)
    (5,19.8)
    (10,20.7)
    (15,20.95)
};
\addlegendentry{$f_1$ und $f_2$}


% ------------------------------------------------------------
% Ideale Grundsignal-Gerade, extrapoliert
% ------------------------------------------------------------
\addplot[
    DARCblue,
    very thick,
    dashed,
    domain=-50:\IIP,
    samples=2
] {x+\G};
\addlegendentry{}


% ------------------------------------------------------------
% Reale IM3-Produkte mit Kompression und gleicher Sättigung
% ------------------------------------------------------------
\addplot[
    DARCred,
    very thick,
    smooth
] coordinates {
    (-30,-90)
    (-25,-75)
    (-20,-60)
    (-15,-45)
    (-10,-30)
    (-5,-15)
    (0,0)
    (3,8.5)
    (5,13.0)
    (7,16.2)
    (9,18.4)
    (11,19.7)
    (13,20.45)
    (15,20.85)
};
\addlegendentry{3. Ordnung}


% ------------------------------------------------------------
% IM3-Produkte: ideale Gerade, extrapoliert
% ------------------------------------------------------------
\addplot[
    DARCred,
    very thick,
    dashed,
    domain=-30:\IIP,
    samples=2
] {3*x};
\addlegendentry{}

% ------------------------------------------------------------
% IP3: Schnittpunkt der extrapolierten Geraden
% ------------------------------------------------------------
\addplot[
    black,
    only marks,
    mark=*,
    mark size=2.4pt
] coordinates {(\IIP,\OIP)};

\draw[densely dashed]
    (axis cs:\IIP,-90) -- (axis cs:\IIP,\OIP);

\draw[densely dashed]
    (axis cs:-50,\OIP) -- (axis cs:\IIP,\OIP);

\node[anchor=south west]
    at (axis cs:\IIP,\OIP)
    {IP3};

\end{axis}

\end{tikzpicture}